% Ad Atticum 12.15 --- headnote
% ad-atticum-12-15, 9 March 45 BC, Astura

\textit{Cicero to Atticus, written from Astura on the
seventh day before the Ides of March 709 AUC --- 9
March 45 BC (the manuscript dateline: \textit{Scr.\
Asturae vii Id.\ Mart.\ a.\ 709 (45)}). A single
section, the most often quoted of the early Astura
letters. After a one-sentence instruction about the
augur Appuleius (the standing excuse will not do, so
the daily excuse must continue), Cicero gives the
sequence's clearest description of his routine:
\textit{cumque mane me in silvam abstrusi densam et
asperam, non exeo inde ante vesperum} --- in the
morning he hides himself in a thick rough wood and
does not come out until evening.}

\textit{The middle sentence --- \textit{secundum te
nihil est mihi amicius solitudine}, ``next to you,
nothing is more my friend than solitude'' --- carries
the keynote that runs through the sequence. In the
wood his talk is with books, but weeping breaks in on
it and he is, so far, no match for it. He closes by
promising the reply to Brutus by tomorrow, to be
forwarded by Atticus when a courier appears.}
